\section{Estado del Arte}

La enseñanza de compiladores e intérpretes continúa siendo uno de los retos más complejos dentro de los programas de ingeniería y ciencias de la computación, debido al nivel de abstracción involucrado en conceptos como lenguajes formales, análisis léxico–sintáctico, representaciones intermedias, optimización y generación de código. Estas dificultades han motivado el desarrollo de herramientas didácticas que integran visualización, simulación e interactividad con el fin de mejorar la comprensión conceptual y reducir la brecha entre teoría e implementación práctica.

Los primeros esfuerzos en esta línea incorporaron mecanismos de visualización para representar las transformaciones internas del proceso de compilación. Vegdahl \cite{Vegdahl2000UsingVT} propuso paquetes gráficos para inspeccionar distintas etapas del compilador, permitiendo observar árboles sintácticos y representaciones intermedias. Posteriormente, Baldwin \cite{Baldwin2003ACompiler} presentó un compilador modular de pequeña escala diseñado específicamente con fines educativos, reduciendo la complejidad de implementación y facilitando la comprensión progresiva de los componentes fundamentales. De manera complementaria, Mernik y Žumer desarrollaron LISA, un entorno integrado capaz de generar intérpretes y compiladores a partir de especificaciones formales, incluyendo animaciones y visualización de autómatas y estructuras semánticas \cite{Mernik2003AnET}.

En años recientes, la investigación ha evolucionado hacia plataformas web interactivas que permiten simulaciones paso a paso y retroalimentación inmediata. ComVIS ofrece un entorno de simulación para autómatas y técnicas de parsing LL/LR, reportando mejoras en la comprensión conceptual y en el desempeño académico de los estudiantes \cite{jovanovic2021comvis}. De forma similar, sistemas web educativos integrales han demostrado que la combinación de visualización dinámica, ejercicios guiados y experimentación directa favorece el aprendizaje de la construcción de compiladores \cite{Stamenković2024AWE}. Los tutores interactivos, como ITT, refuerzan este enfoque al conectar la implementación real con explicaciones y seguimiento guiado, mostrando incrementos en la retención de contenidos teóricos \cite{delvado2023itt}.

Otra línea de trabajo se centra en herramientas de inspección y exploración del comportamiento interno del compilador. Propuestas como Compiler Explorer permiten analizar en tiempo real el código intermedio y el ensamblador generado, facilitando la comprensión de optimizaciones y decisiones del compilador \cite{compilerexplorer2022}. Asimismo, soluciones recientes como VisOpt incorporan visualizaciones específicas de procesos de optimización para apoyar la enseñanza de etapas avanzadas del pipeline de compilación \cite{koitzhristov2025visopt}.

El uso de proyectos didácticos y compiladores educativos autocontenidos también ha sido ampliamente explorado. Estudios reportan que la construcción de compiladores pequeños o modulares incrementa la participación estudiantil y favorece la transferencia de conocimientos prácticos \cite{Kundra2016AnER,liu2017evaluation}. En este contexto, se han propuesto diseños modulares con comunicación entre fases \cite{navas2022modular}, reimplementaciones simplificadas de compiladores completos \cite{berezun2022reimplementing} y herramientas pedagógicas para la visualización de intérpretes \cite{Spigariol2023Aprendiendo}.

Desde una perspectiva comparativa, investigaciones empíricas han evaluado tecnologías de generación de analizadores sintácticos. Ortín et al.\ \cite{ortin2022antlr} compararon Lex/Yacc y ANTLR en contextos educativos, evidenciando mayor facilidad de uso, mantenibilidad y mejores tasas de aprobación a favor de ANTLR. Adicionalmente, marcos modernos orientados a claridad semántica y corrección formal, como ChocoPy y CakeML, han sido analizados por su potencial valor didáctico al proporcionar implementaciones más comprensibles y verificables \cite{daleman2024chocopy,myreen2021cakeml}.

Finalmente, revisiones sistemáticas en tecnología educativa destacan que la efectividad de las herramientas de aprendizaje depende no solo de su funcionalidad técnica, sino también de aspectos como usabilidad, retroalimentación inmediata y soporte a la experimentación activa \cite{lu2022usability}. Estos criterios se han convertido en factores clave para el diseño de entornos educativos modernos.

En síntesis, el estado del arte evidencia una tendencia hacia entornos interactivos, visuales y modulares que integran simulación, experimentación práctica y evaluación empírica del aprendizaje. No obstante, persisten desafíos relacionados con la escalabilidad, la integración curricular y la validación sistemática del impacto pedagógico, lo que justifica el desarrollo de nuevas propuestas que combinen simplicidad conceptual, interactividad y soporte experimental.
